\begin{landscape}
\thispagestyle{empty}

\section*{ANEXOS}
\addcontentsline{toc}{section}{ANEXOS}

\begin{center}
    \textbf{\Large MATRIZ DE CONSISTENCIA} \\
    \vspace{0.2cm}
    \small Título: Efecto de un sistema de gestión local con integridad de datos y facturación electrónica en la optimización operativa de los comercios de abarrotes, Cusco 2026
\end{center}

\vspace{0.3cm}

\noindent
\begin{tabularx}{\linewidth}{|Y|Y|Y|Y|Y|}
\hline
\rowcolor[gray]{0.9} 
\textbf{Problemas} & \textbf{Objetivos} & \textbf{Hipótesis} & \textbf{Variables} & \textbf{Metodología} \\ \hline

\textbf{Problema General:} \newline
¿Qué efecto tiene la implementación de un sistema de gestión local con integridad de datos y facturación electrónica directa en la optimización de los procesos operativos en comercios de abarrotes de Cusco, 2026? \newline

\textbf{Problemas Específicos:}
\begin{itemize}[leftmargin=*, nosep, itemsep=6pt]
    \item \textbf{PE1:} ¿En qué medida la aplicación de un motor de validación transaccional reduce la incidencia de errores de inventario?
    \item \textbf{PE2:} ¿De qué manera la integración directa de facturación bajo estándar UBL influye en el tiempo de atención al cliente?
    \item \textbf{PE3:} ¿Cómo influye la arquitectura de datos local en la confiabilidad de los reportes de gestión?
\end{itemize} 
& 
\textbf{Objetivo General:} \newline
Evaluar el efecto de la implementación de un sistema de gestión local con integridad de datos y facturación electrónica en la optimización de los procesos operativos en comercios de abarrotes. \newline

\textbf{Objetivos Específicos:}
\begin{itemize}[leftmargin=*, nosep, itemsep=6pt]
    \item \textbf{OE1:} Comprobar la reducción de errores de inventario mediante la aplicación de un motor de validación.
    \item \textbf{OE2:} Determinar el impacto de la facturación electrónica directa en la disminución del tiempo de emisión de comprobantes.
    \item \textbf{OE3:} Evaluar la mejora en la confiabilidad de los reportes analíticos a partir de una arquitectura de datos local.
\end{itemize} 
& 
\textbf{Hipótesis General:} \newline
La implementación de un sistema de gestión local con integridad de datos y facturación electrónica directa produce una optimización significativa en los procesos operativos de los comercios de abarrotes. \newline

\textbf{Hipótesis Específicas:}
\begin{itemize}[leftmargin=*, nosep, itemsep=6pt]
    \item \textbf{HIP1:} El motor de validación transaccional reduce significativamente la incidencia de stock negativo e inconsistente.
    \item \textbf{HIP2:} La integración directa con SUNAT disminuye el tiempo promedio de facturación frente a soluciones externas.
    \item \textbf{HIP3:} La arquitectura local garantiza mayor disponibilidad y veracidad de la información para la toma de decisiones.
\end{itemize} 
& 
\textbf{Variable Independiente (X):}
\begin{itemize}[leftmargin=*, nosep, itemsep=2pt]
    \item Sistema de gestión local con integridad y facturación directa.
\end{itemize}

\textbf{Variables Dependientes (Y):}
\begin{itemize}[leftmargin=*, nosep, itemsep=2pt]
    \item Eficiencia en el control de inventario (Errores).
    \item Tiempo de atención (Emisión de CPE).
    \item Confiabilidad de reportes analíticos.
\end{itemize} 
& 
\textbf{Nivel:} Explicativo (Enfoque Experimental). \newline
\textbf{Método:} Cuasi-experimental[cite: 100]. \newline
\textbf{Diseño:} Pre-experimental con pre-prueba y post-prueba[cite: 106]. \newline

\textbf{Muestra:} \newline
500 operaciones de venta y 50 categorías de productos[cite: 120]. \newline

\textbf{Técnicas:} Experimentación controlada y observación[cite: 132]. \newline
\textbf{Instrumentos:} Suite de pruebas automatizadas (Pytest) y Guía de observación técnica[cite: 115, 133]. \\ \hline
\end{tabularx}
\end{landscape}